\documentclass[a4paper,12pt]{article}

% -------------------------------------------------
% Кодировка, язык, математика
% -------------------------------------------------
\usepackage[T2A]{fontenc}
\usepackage[utf8]{inputenc}
\usepackage[russian]{babel}

\usepackage{amsmath,amssymb,amsthm,mathtools}
\usepackage{icomma}

% -------------------------------------------------
% Типографика
% -------------------------------------------------
\usepackage{helvet}
\renewcommand{\familydefault}{\sfdefault}

\usepackage{geometry}
\geometry{
  left=18mm,
  right=18mm,
  top=20mm,
  bottom=20mm
}

\usepackage{microtype}
\usepackage{parskip}
\usepackage{setspace}
\onehalfspacing

% -------------------------------------------------
% Графика
% -------------------------------------------------
\usepackage{tikz}
\usepackage{tkz-euclide}
\usepackage{pgfplots}
\pgfplotsset{compat=1.18}

% -------------------------------------------------
% Таблицы и списки
% -------------------------------------------------
\usepackage{tabularx,booktabs,longtable}
\usepackage{enumitem}
\setlist[itemize]{
  leftmargin=6mm,
  itemsep=2pt,
  topsep=3pt
}
\setlist[enumerate]{
  leftmargin=8mm,
  itemsep=5pt,
  topsep=4pt
}

% -------------------------------------------------
% Цвет и ссылки
% -------------------------------------------------
\usepackage{xcolor}

\definecolor{accent}{HTML}{008080}
\definecolor{darkaccent}{HTML}{004D4D}
\definecolor{textgray}{HTML}{333333}
\definecolor{softgray}{HTML}{6A7373}
\definecolor{linegray}{HTML}{D7DEDE}

\usepackage{hyperref}
\hypersetup{
  colorlinks=true,
  linkcolor=darkaccent,
  urlcolor=darkaccent
}

% -------------------------------------------------
% Колонтитулы
% -------------------------------------------------
\usepackage{fancyhdr}
\pagestyle{fancy}
\fancyhf{}
\fancyhead[L]{\small\textcolor{softgray}{Повторение: геометрия, векторы, стереометрия}}
\fancyhead[R]{\small\textcolor{softgray}{Тимофей Васильченко}}
\fancyfoot[C]{\small\textcolor{softgray}{\thepage}}
\renewcommand{\headrulewidth}{0.3pt}
\renewcommand{\headrule}{%
  \hbox to\headwidth{\color{linegray}\leaders\hrule height \headrulewidth\hfill}
}

% -------------------------------------------------
% Служебные команды
% -------------------------------------------------
\setlength{\parindent}{0pt}
\setlength{\emergencystretch}{2em}

\newcommand{\topic}[1]{%
  \vspace{0.8em}
  {\Large\bfseries\textcolor{darkaccent}{#1}}\par
  \vspace{0.25em}
  {\color{accent}\rule{\linewidth}{0.6pt}}
  \vspace{0.45em}
}

\newcommand{\subtopic}[1]{%
  \vspace{0.45em}
  {\large\bfseries\textcolor{accent}{#1}}\par
  \vspace{0.2em}
}

\newcommand{\key}[1]{\textcolor{darkaccent}{\textbf{#1}}}
\newcommand{\vect}[1]{\overrightarrow{#1}}
\newcommand{\R}{\mathbb{R}}

\begin{document}

% =================================================
% ТИТУЛЬНЫЙ ЛИСТ
% =================================================
\begin{titlepage}
\thispagestyle{empty}

\begin{center}

\vspace*{0.19\textheight}

{\Huge\bfseries\textcolor{darkaccent}{Чек-лист}}

\vspace{8mm}

{\LARGE\bfseries
Геометрия тестовой части:\\[2mm]
окружность, четырёхугольники,\\
треугольники и трапеции}

\vspace{5mm}

{\Large\bfseries
Векторы и базовая стереометрия}

\vspace{18mm}

{\large
Лёвин Артём Александрович\\[2mm]
\textbf{эксклюзивно для Тимофея Васильченко}
}

\vfill

{\large \today}

\end{center}

% Сдержанная титульная кривая Безье
\begin{tikzpicture}[remember picture,overlay]
\draw[
  accent!55,
  line width=1.2pt
]
([xshift=18mm,yshift=18mm]current page.south west)
.. controls
([xshift=55mm,yshift=28mm]current page.south west)
and
([xshift=-55mm,yshift=8mm]current page.south east)
..
([xshift=-18mm,yshift=18mm]current page.south east);
\end{tikzpicture}

\end{titlepage}

% =================================================
% ОСНОВНОЙ ЧЕК-ЛИСТ
% =================================================

\topic{1. Главный алгоритм геометрической задачи}

Перед вычислениями полезно провести короткое \key{сканирование чертежа}.

\begin{enumerate}
  \item Что известно об углах, сторонах, серединах и параллельных прямых?
  \item Какие \key{треугольники} уже присутствуют на рисунке?
  \item Можно ли получить удобный треугольник дополнительным построением?
  \item Есть ли прямоугольный или равнобедренный треугольник?
  \item Какие углы или отрезки можно сразу признать равными?
  \item Можно ли свести искомую величину к разности, сумме или отношению уже знакомых величин?
\end{enumerate}

\[
\boxed{\text{Сложная фигура}
\;\longrightarrow\;
\text{удобные треугольники}
\;\longrightarrow\;
\text{известные свойства}}
\]

\key{Практический принцип.}
Если решение не видно, сначала ищите не новую формулу, а удобный треугольник.

% =================================================

\topic{2. Окружность: вписанный угол и дуга}

\subtopic{Базовое свойство}

Если вершина угла лежит на окружности, а его стороны пересекают окружность, то угол является \key{вписанным}.

\[
\boxed{
\angle ABC=
\frac12\,\overset{\frown}{AC}
}
\]

Здесь берётся именно та дуга \(AC\), которая \key{не содержит вершину \(B\)}.

Следовательно,

\[
\boxed{
\overset{\frown}{AC}=2\angle ABC
}
\]

\subtopic{Типовая схема с четырёхугольником}

Пусть \(ABCD\) вписан в окружность,

\[
\angle ABC=103^\circ,
\qquad
\angle CAD=42^\circ.
\]

Требуется найти \(\angle ABD\).

Угол \(\angle ABC\) опирается на дугу \(ADC\), поэтому

\[
\overset{\frown}{ADC}
=
2\cdot103^\circ
=
206^\circ.
\]

Угол \(\angle CAD\) опирается на дугу \(CD\), значит

\[
\overset{\frown}{CD}
=
2\cdot42^\circ
=
84^\circ.
\]

Отсюда

\[
\overset{\frown}{AD}
=
206^\circ-84^\circ
=
122^\circ.
\]

Следовательно,

\[
\boxed{
\angle ABD
=
\frac12\cdot122^\circ
=
61^\circ
}
\]

\subtopic{Что проверять}

\begin{itemize}
  \item Не перепутана ли малая дуга с большой?
  \item Не содержит ли выбранная дуга вершину вписанного угла?
  \item При переходе от угла к дуге нужно \key{умножать на \(2\)}.
  \item При переходе от дуги к углу нужно \key{делить на \(2\)}.
\end{itemize}

% =================================================

\topic{3. Площадь параллелограмма и середина стороны}

Пусть \(ABCD\) --- параллелограмм, \(E\) --- середина \(AD\).

\begin{center}
\begin{tikzpicture}[scale=0.9]
  \tkzDefPoint(0,0){A}
  \tkzDefPoint(4.8,0){B}
  \tkzDefPoint(6.2,2.8){C}
  \tkzDefPoint(1.4,2.8){D}
  \tkzDefPoint(0.7,1.4){E}

  \tkzDrawPolygon[thick,color=darkaccent](A,B,C,D)
  \tkzDrawSegment[thick,color=accent](B,E)
  \tkzDrawPoints[color=darkaccent,fill=darkaccent,size=3](A,B,C,D,E)

  \tkzLabelPoints[below](A,B)
  \tkzLabelPoints[above](C,D)
  \tkzLabelPoints[left](E)
\end{tikzpicture}
\end{center}

Если

\[
S_{ABCD}=S,
\]

то

\[
AE=\frac{AD}{2}.
\]

Треугольник \(ABE\) и параллелограмм \(ABCD\) имеют одну и ту же высоту \(h\), проведённую к прямой \(AD\).

Поэтому

\[
S_{ABE}
=
\frac12\cdot AE\cdot h
=
\frac12\cdot\frac{AD}{2}\cdot h
=
\frac14 ADh.
\]

Но

\[
ADh=S_{ABCD}.
\]

Следовательно,

\[
\boxed{
S_{ABE}=\frac14S_{ABCD}
}
\]

и

\[
\boxed{
S_{BCDE}=\frac34S_{ABCD}.
}
\]

Если \(S_{ABCD}=24\), то

\[
S_{ABE}=6,
\qquad
S_{BCDE}=24-6=18.
\]

\key{Надёжный метод:}
не полагайтесь на визуальное ощущение площади. Связывайте площади через основания и общую высоту.

% =================================================

\topic{4. Прямоугольный треугольник: медиана к гипотенузе}

Пусть

\[
\angle ACB=90^\circ,
\]

а \(M\) --- середина гипотенузы \(AB\).

Тогда

\[
\boxed{
MA=MB=MC=\frac{AB}{2}
}
\]

То есть середина гипотенузы равноудалена от всех трёх вершин.

\begin{center}
\begin{tikzpicture}[scale=1.05]
  \tkzDefPoint(4,0){A}
  \tkzDefPoint(0,3){B}
  \tkzDefPoint(0,0){C}
  \tkzDefPoint(2,1.5){M}
  \tkzDefPoint(1.44,1.92){H}

  \tkzDrawPolygon[thick,color=darkaccent](A,B,C)
  \tkzDrawSegment[thick,color=accent](C,M)
  \tkzDrawSegment[thick,color=accent](C,H)

  \tkzDrawPoints[color=darkaccent,fill=darkaccent,size=3](A,B,C,M,H)

  \tkzMarkRightAngle[size=0.28,color=darkaccent](A,C,B)
  \tkzMarkRightAngle[size=0.22,color=accent](C,H,B)

  \tkzLabelPoints[below](A,C)
  \tkzLabelPoints[above](B)
  \tkzLabelPoints[right](M)
  \tkzLabelPoints[left](H)
\end{tikzpicture}
\end{center}

Из

\[
MB=MC
\]

следует, что \(\triangle BMC\) равнобедренный, поэтому

\[
\boxed{
\angle MBC=\angle BCM
}
\]

Но луч \(BM\) лежит на гипотенузе \(BA\), следовательно,

\[
\angle MBC=\angle ABC.
\]

Итак,

\[
\boxed{
\angle BCM=\angle ABC.
}
\]

\subtopic{Медиана и высота из прямого угла}

Пусть дополнительно

\[
CH\perp AB.
\]

Если

\[
\angle ABC=\alpha,
\]

то в прямоугольном треугольнике \(BCH\)

\[
\angle BCH=90^\circ-\alpha.
\]

При этом

\[
\angle BCM=\alpha.
\]

Поэтому угол между медианой \(CM\) и высотой \(CH\) равен

\[
\boxed{
\angle HCM
=
\left|
\alpha-(90^\circ-\alpha)
\right|
=
|2\alpha-90^\circ|.
}
\]

Например, при

\[
\alpha=65^\circ
\]

имеем

\[
\angle BCH=25^\circ,
\qquad
\angle BCM=65^\circ,
\]

следовательно,

\[
\boxed{
\angle HCM=40^\circ.
}
\]

\subtopic{Важная проверка чертежа}

Если при вычислении обычного геометрического угла получилось отрицательное число, нельзя просто менять знак.

Нужно проверить:

\begin{itemize}
  \item в каком порядке реально расположены лучи;
  \item какой угол является частью другого;
  \item корректно ли построены медиана, высота и биссектриса.
\end{itemize}

Для угла между двумя лучами удобно использовать модуль разности их направлений.

% =================================================

\topic{5. Биссектриса}

Биссектриса делит угол на две равные части.

Если \(BK\) --- биссектриса \(\angle ABC\), то

\[
\boxed{
\angle ABK=\angle KBC
=
\frac12\angle ABC.
}
\]

В прямоугольном треугольнике при построении одновременно медианы, высоты и биссектрисы важно сначала определить \key{порядок лучей}, а уже затем вычитать углы.

Например, если прямой угол равен \(90^\circ\), его биссектриса образует с каждой стороной угол

\[
45^\circ.
\]

% =================================================

\topic{6. Трапеция: средняя линия и диагональ}

Пусть \(ABCD\) --- трапеция,

\[
AD\parallel BC,
\]

\(E\) и \(F\) --- середины боковых сторон \(AB\) и \(CD\).

Средняя линия трапеции равна полусумме оснований:

\[
\boxed{
EF=\frac{AD+BC}{2}.
}
\]

Пусть диагональ \(BD\) пересекает среднюю линию в точке \(O\).

В треугольнике \(ABD\):

\[
E \text{ --- середина } AB.
\]

Поскольку \(EO\parallel AD\), отрезок \(EO\) является средней линией треугольника \(ABD\). Поэтому

\[
\boxed{
EO=\frac{AD}{2}.
}
\]

Аналогично в треугольнике \(BCD\)

\[
\boxed{
OF=\frac{BC}{2}.
}
\]

Следовательно, диагональ делит среднюю линию трапеции на отрезки, равные \key{половинам оснований}.

Если основания равны \(10\) и \(4\), то получаются отрезки

\[
5
\quad\text{и}\quad
2.
\]

Больший из них:

\[
\boxed{5}.
\]

\key{Ключевой ход:}
рассматривайте не всю трапецию, а два треугольника, возникающих после проведения диагонали.

% =================================================

\topic{7. Векторы на координатной плоскости}

\subtopic{Координаты вектора}

Если

\[
A(x_A;y_A),
\qquad
B(x_B;y_B),
\]

то

\[
\boxed{
\vect{AB}
=
(x_B-x_A;\,y_B-y_A).
}
\]

На клетчатой плоскости координаты можно также читать как перемещения:

\begin{itemize}
  \item вправо --- \(+\) по \(x\);
  \item влево --- \(-\) по \(x\);
  \item вверх --- \(+\) по \(y\);
  \item вниз --- \(-\) по \(y\).
\end{itemize}

\subtopic{Длина вектора}

Для

\[
\vec a=(a_x;a_y)
\]

длина равна

\[
\boxed{
|\vec a|
=
\sqrt{a_x^2+a_y^2}.
}
\]

\subtopic{Умножение на число}

Если

\[
\vec a=(a_x;a_y),
\]

то

\[
\boxed{
k\vec a=(ka_x;ka_y).
}
\]

\subtopic{Сложение и вычитание}

\[
\boxed{
\vec a+\vec b
=
(a_x+b_x;\,a_y+b_y)
}
\]

\[
\boxed{
\vec a-\vec b
=
(a_x-b_x;\,a_y-b_y)
}
\]

\subtopic{Скалярное произведение}

Через координаты:

\[
\boxed{
\vec a\cdot\vec b
=
a_xb_x+a_yb_y
}
\]

Через длины и угол \(\varphi\) между векторами:

\[
\boxed{
\vec a\cdot\vec b
=
|\vec a|\,|\vec b|\cos\varphi.
}
\]

Отсюда, если оба вектора ненулевые,

\[
\boxed{
\cos\varphi
=
\frac{\vec a\cdot\vec b}
{|\vec a|\,|\vec b|}.
}
\]

\subtopic{Пример 1}

Пусть

\[
\vec a=(4;6),
\qquad
\vec b=(6;-2).
\]

Тогда

\[
\vec a\cdot\vec b
=
4\cdot6+6\cdot(-2)
=
24-12
=
\boxed{12}.
\]

\key{Типичная ошибка:}
потерять минус у координаты вектора, направленного вниз или влево.

\subtopic{Пример 2}

Пусть

\[
\vec a=(2;0),
\qquad
\vec b=(1;4).
\]

Требуется найти

\[
|\vec a+3\vec b|.
\]

Сначала

\[
3\vec b=(3;12).
\]

Тогда

\[
\vec a+3\vec b
=
(2;0)+(3;12)
=
(5;12).
\]

Следовательно,

\[
|\vec a+3\vec b|
=
\sqrt{5^2+12^2}
=
\sqrt{169}
=
\boxed{13}.
\]

\key{Экзаменационный алгоритм:}

\[
\boxed{
\text{сначала координаты}
\;\longrightarrow\;
\text{потом длина или скалярное произведение}
}
\]

% =================================================

\topic{8. Цилиндр: отношения размеров и объёмов}

Объём цилиндра:

\[
\boxed{
V=\pi r^2h.
}
\]

Здесь

\[
r \text{ --- радиус},
\qquad
h \text{ --- высота}.
\]

\subtopic{«В \(k\) раз шире»}

Для цилиндра ширину определяет радиус или диаметр.

Если второй цилиндр в \(k\) раз шире первого, то

\[
r_2=kr_1.
\]

Так как радиус в формуле объёма стоит \key{в квадрате},

\[
r_2^2=k^2r_1^2.
\]

Поэтому увеличение радиуса в \(2\) раза увеличивает площадь основания и объём при неизменной высоте в

\[
2^2=4
\]

раза.

\subtopic{Отношение объёмов}

\[
\frac{V_2}{V_1}
=
\frac{\pi r_2^2h_2}
{\pi r_1^2h_1}
=
\boxed{
\frac{r_2^2h_2}{r_1^2h_1}
}.
\]

Если

\[
h_1=mh_2,
\qquad
r_2=kr_1,
\]

то

\[
\boxed{
\frac{V_2}{V_1}=\frac{k^2}{m}.
}
\]

Например,

\[
h_1=2h_2,
\qquad
r_2=1{,}5r_1.
\]

Тогда

\[
\frac{V_2}{V_1}
=
\frac{(1{,}5)^2}{2}
=
\frac{2{,}25}{2}
=
1{,}125
=
\boxed{\frac98}.
\]

\subtopic{Переливание жидкости}

При переливании без потерь объём жидкости \key{не меняется}:

\[
V_1=V_2.
\]

Следовательно,

\[
\pi r_1^2h_1
=
\pi r_2^2h_2.
\]

Если

\[
r_2=kr_1,
\]

то

\[
r_1^2h_1
=
k^2r_1^2h_2,
\]

откуда

\[
\boxed{
h_2=\frac{h_1}{k^2}.
}
\]

Если

\[
h_1=16,
\qquad
d_2=2d_1,
\]

то

\[
r_2=2r_1
\]

и

\[
h_2
=
\frac{16}{4}
=
\boxed{4}.
\]

\key{Важно:}

\[
\frac{d_2}{d_1}
=
\frac{r_2}{r_1},
\]

поэтому при сравнении диаметров можно сразу перейти к тому же отношению радиусов.

% =================================================

\topic{9. Пирамида внутри прямоугольного параллелепипеда}

Объём пирамиды:

\[
\boxed{
V_{\text{пир}}
=
\frac13S_{\text{осн}}h.
}
\]

Если основание пирамиды --- прямоугольник со сторонами \(a\) и \(b\), то

\[
S_{\text{осн}}=ab.
\]

Поэтому

\[
\boxed{
V_{\text{пир}}
=
\frac13abh.
}
\]

Для

\[
a=3,
\qquad
b=9,
\qquad
h=4
\]

получаем

\[
V_{\text{пир}}
=
\frac13\cdot3\cdot9\cdot4
=
\boxed{36}.
\]

Объём прямоугольного параллелепипеда с теми же основанием и высотой:

\[
V_{\text{пар}}=abh.
\]

Следовательно,

\[
\boxed{
V_{\text{пир}}
=
\frac13V_{\text{пар}}.
}
\]

Это позволяет решать задачи даже тогда, когда длины рёбер не даны, но известен объём параллелепипеда.

% =================================================

\topic{10. Типичные ошибки и финальная самопроверка}

\begin{tabularx}{\linewidth}{@{}p{0.30\linewidth}X@{}}
\toprule
\textbf{Ситуация} & \textbf{Что проверить} \\
\midrule

Вписанный угол
&
Правильно ли выбрана дуга; не содержит ли она вершину угла; не перепутаны ли множители \(2\) и \(\frac12\).
\\[2mm]

Площади
&
Используются ли действительно общие высоты; не сделан ли вывод только «по рисунку».
\\[2mm]

Прямоугольный треугольник
&
Сумма острых углов равна \(90^\circ\); медиана к гипотенузе равна её половине.
\\[2mm]

Несколько лучей из одной вершины
&
Определён ли их реальный порядок; корректно ли выполнено вычитание углов.
\\[2mm]

Трапеция
&
Не забыто ли, что средняя линия равна полусумме оснований, а диагональ делит её на половины оснований.
\\[2mm]

Векторы
&
Сохранены ли знаки координат; сначала ли найдены координаты результирующего вектора.
\\[2mm]

Цилиндр
&
Не забыта ли вторая степень радиуса; правильно ли составлено требуемое отношение \(V_2/V_1\).
\\[2mm]

Переливание
&
Записано ли равенство объёмов \(V_1=V_2\).
\\[2mm]

Пирамида
&
Не потерян ли коэффициент \(\frac13\).
\\

\bottomrule
\end{tabularx}

\vspace{6mm}

\key{Короткая проверка перед ответом:}

\[
\boxed{
\text{условие}
\to
\text{чертёж}
\to
\text{нужная фигура}
\to
\text{формула}
\to
\text{подстановка}
\to
\text{арифметика}
\to
\text{проверка}
}
\]

% =================================================
% ТРЕНИРОВОЧНЫЙ БЛОК
% =================================================

\clearpage

\topic{Тренировочный блок}

\textbf{Упражнения 1--3 --- средний уровень.}

\begin{enumerate}

\item
Четырёхугольник \(ABCD\) вписан в окружность.

Известно, что

\[
\angle ABC=108^\circ,
\qquad
\angle CAD=37^\circ.
\]

Найдите \(\angle ABD\).

\item
Площадь параллелограмма \(ABCD\) равна \(40\).
Точка \(E\) --- середина стороны \(AD\).

Найдите площадь четырёхугольника \(BCDE\).

\item
В прямоугольном треугольнике \(ABC\)

\[
\angle C=90^\circ,
\qquad
\angle B=58^\circ.
\]

Из вершины \(C\) проведены медиана \(CM\) к гипотенузе и высота \(CH\).

Найдите \(\angle HCM\).

\end{enumerate}

\vspace{3mm}

\textbf{Упражнения 4--9 --- уровень выше среднего.}

\begin{enumerate}[resume]

\item
Основания трапеции равны \(6\) и \(14\).
Одна из диагоналей делит среднюю линию трапеции на два отрезка.

Найдите больший из них.

\item
Даны векторы

\[
\vec a=(3;-2),
\qquad
\vec b=(-1;5).
\]

Найдите

\[
\vec a\cdot\vec b.
\]

\item
Даны векторы

\[
\vec a=(1;2),
\qquad
\vec b=(2;-1).
\]

Найдите

\[
|\vec a+2\vec b|.
\]

\item
Первая цилиндрическая кружка в \(3\) раза выше второй, а радиус второй в \(2\) раза больше радиуса первой.

Найдите

\[
\frac{V_2}{V_1}.
\]

\item
Уровень жидкости в первом цилиндрическом сосуде равен \(18\) см.
Жидкость полностью перелили во второй цилиндрический сосуд, диаметр которого в \(3\) раза больше диаметра первого.

Найдите высоту уровня жидкости во втором сосуде.

\item
Основание прямоугольного параллелепипеда имеет размеры

\[
5\times8,
\]

а его высота равна \(6\).

Найдите объём пирамиды, основание которой совпадает с основанием параллелепипеда, а вершина лежит в плоскости его верхнего основания.

\end{enumerate}

\vspace{3mm}

\textbf{Упражнение 10 --- сложное.}

\begin{enumerate}[resume]

\item
В прямоугольном треугольнике \(ABC\)

\[
\angle C=90^\circ,
\qquad
\angle B=67^\circ.
\]

Из вершины \(C\) проведены:

\begin{itemize}
  \item высота \(CH\) к гипотенузе;
  \item биссектриса \(CL\);
  \item медиана \(CM\) к гипотенузе.
\end{itemize}

Найдите

\[
\angle HCL,
\qquad
\angle LCM,
\qquad
\angle HCM.
\]

\end{enumerate}

% =================================================
% ОТВЕТЫ
% =================================================

\clearpage

\topic{Ответы к тренировочному блоку}

\renewcommand{\arraystretch}{1.35}

\begin{table}[ht]
\centering
\caption{Ответы}
\begin{tabularx}{\linewidth}{@{}c p{0.22\linewidth} X@{}}
\toprule
\textbf{№} & \textbf{Ответ} & \textbf{Ключевая идея для самопроверки} \\
\midrule

1
&
\(71^\circ\)
&
Дуги \(ADC\) и \(CD\): \(216^\circ\) и \(74^\circ\); дуга \(AD=142^\circ\).
\\

2
&
\(30\)
&
Треугольник \(ABE\) занимает \(\frac14\) площади параллелограмма.
\\

3
&
\(26^\circ\)
&
\(\angle BCM=58^\circ\), \(\angle BCH=32^\circ\).
\\

4
&
\(7\)
&
Диагональ делит среднюю линию на половины оснований: \(3\) и \(7\).
\\

5
&
\(-13\)
&
\(3\cdot(-1)+(-2)\cdot5=-13\).
\\

6
&
\(5\)
&
\(\vec a+2\vec b=(5;0)\).
\\

7
&
\(\dfrac43\)
&
\(\dfrac{V_2}{V_1}=\dfrac{2^2}{3}\).
\\

8
&
\(2\text{ см}\)
&
Радиус увеличился в \(3\) раза, площадь основания --- в \(9\) раз.
\\

9
&
\(80\)
&
\(\frac13\cdot5\cdot8\cdot6=80\).
\\

10
&
\(22^\circ,\ 22^\circ,\ 44^\circ\)
&
От луча \(CB\): высота --- \(23^\circ\), биссектриса --- \(45^\circ\), медиана --- \(67^\circ\).
\\

\bottomrule
\end{tabularx}
\end{table}

\vfill

\begin{center}
{\small\textcolor{softgray}{
Лёвин Артём Александрович эксклюзивно для Тимофея Васильченко
}}
\end{center}

\end{document}